\documentclass[a4paper, 10pt]{article}
\usepackage{palatino}
\usepackage[a4paper]{geometry}
\usepackage[UKenglish]{babel}
\usepackage[utf8]{inputenc}
\usepackage{afterpage}
\usepackage{fancyhdr}
\usepackage{amsmath}
\usepackage{amstext}
\usepackage{amssymb}
\usepackage{physics}
\usepackage{derivative}
\usepackage{longtable}
\usepackage{tabularx}
\usepackage{dcolumn}
\usepackage{graphicx}
\usepackage{enumerate}
\usepackage{multicol}
\usepackage{pdfpages}
\usepackage{lastpage}
\usepackage{siunitx}
\usepackage{listings}
\usepackage{tikz}
\usepackage{xspace}
\usepackage{caption} 

\AtBeginDocument{\RenewCommandCopy\qty\SI}
\sloppy

\geometry{margin=25mm, top=20mm, bottom=20mm}
\captionsetup{format=plain,indention=0mm, font=small, labelfont=bf, labelsep=period}

\newcommand{\logo}{\begin{figure}[!t] \vspace*{-10mm} \hspace*{-10mm}
\includegraphics[height=23mm]{../defs/usn-logo-en.pdf} \end{figure}}

\newcommand{\labfoot}[1]{\fancyhf[FC]{ \textit{TSE2280, Lab exercise  #1, spring 2024} } }

\newcommand{\regneovinglosningfoot}[1]{\fancyhf[FC]{TSE2280 Regneøving #1 - Løsning }}
\newcommand{\regneovingtitle}[1]{\section*{TSE2280, våren 2024. \qquad Regneøving #1 }}
\newcommand{\regneovinglosningtitle}[1]{\section*{TSE2280, Måleteknikk og signalbehandling \\Løsning regneøving #1 }}

\newcommand{\Fig}[1]{Fig.~\ref{fig:#1}}
\newcommand{\tableref}[1]{Table~\ref{tab:#1}}

\newcommand{\ed}[1]{\ensuremath{\cdot10^{#1}}}
\renewcommand{\deg}{\ensuremath{^\circ}}
\newcommand{\omegahat}{\hat \omega}

\newcommand{\zm}[1]{ z^{-{#1} } }

\newcommand{\ew}[1]{e^{{#1}j \omegahat }}
\newcommand{\ewn}[2]{e^{{#1}j{#2}} }
\newcommand{\Hw}[1]{H\left( e^{{#1}j \omegahat} \right)}
\newcommand{\Hwn}[2]{H \left( e^{{#1}j {#2}} \right)}

\newcounter{problem}
\newcommand{\problem}[1]{\stepcounter{problem}\section*{Problem \arabic{problem}. {#1}}} 

\newcommand{\subproblem}{\begin{enumerate}[a)]}
\newcommand{\subproblemend}{\end{enumerate}}

%=== FIGURE DEFINITIONS ============
\newenvironment{stdfig}{\begin{figure}[!b]\begin{center}}{\end{center}\end{figure}}
\newcommand{\stdgraphics}[1]{\begin{center}\includegraphics[width=0.95\textwidth]{#1}\end{center}}
\newcommand{\scaledgraphics}[2]{\begin{center}\includegraphics[width={#1}\textwidth]{#2}\end{center}}
\newcommand{\includescanFB}[1]{\begin{center}\includegraphics[width=0.90\textwidth]{#1}\end{center}}
\newcommand{\includescanMcC}[1]{\includegraphics[width=0.80\textwidth]{#1}}
\newcommand{\includedrawing}[1]{\includegraphics[scale=0.65]{#1}}
\newcommand{\Matlabfigure}[1]{\includegraphics[scale=0.80]{#1}}

\newcommand{\losning}[1]{\includegraphics[width=0.85\textwidth]{./losninger/#1}}
\newcommand{\losningfig}[1]{\includegraphics[width=0.85\textwidth]{./figs/#1}}
\newcommand{\losningfigfull}[1]{\includegraphics[width=0.99\textwidth]{./figs/#1}}
\newcommand{\losningfighalf}[1]{\includegraphics[width=0.49\textwidth]{./figs/#1}}
\newcommand{\losningfigreduced}[1]{\includegraphics[width=0.49\textwidth]{./figs/#1}}

%=== REFERENCING ===========================================
\newcommand{\Eq}[1]{(\ref{eq:#1})}

% Headers and footers
\pagestyle{fancy}
\renewcommand{\headrulewidth}{0pt}
\renewcommand{\footrulewidth}{0pt}

\fancyhf{} 
\fancyhf[FL]{\textit{\small LH, \today} }
\fancyhf[FR]{Page \thepage\ of \pageref{LastPage}} 


% Style for program listing
\definecolor{codegreen}{rgb}{0, 0.5, 0}
\definecolor{codegray}{rgb}{0.5, 0.5, 0.5}
\definecolor{codepurple}{rgb}{0.58, 0, 0.82}
\definecolor{backcolour}{rgb}{0.98, 0.98, 0.95}

\lstdefinestyle{pythonstyle}{
	language=python,	
	belowcaptionskip=1\baselineskip,
	breaklines=true,
	frame=single,
	numbers=none, 
	breaklines=true,				 
	keepspaces=true,		
	breakatwhitespace=false,		 		 	
	basicstyle=\small\ttfamily,
	keywordstyle=\color{blue},
	stringstyle=\color{codepurple},	
	commentstyle=\color{codegreen},
	identifierstyle=\color{black!95!white},
	backgroundcolor=\color{backcolour},
	tabsize=4,	
	}
\labfoot{1}

\begin{document}
\logo

\title{Introduction to Complex Exponentials \\ Direction Finding} 
\author{TSE2280 Signal Processing}
\date{Lab 1, Spring 2025}
\maketitle

\thispagestyle{fancy}	

\section{Introduction}
\suppressfloats[t]
\subsection{Background and Motivation}
This lab is a modified version of the lab \emph{Lab P-3: Introduction to Complex Exponentials – Direction Finding}\cite{mcclellan_lab_2016} that accompanies the course text-book \emph{DSP First} by McClellan et~al.~\cite{mcclellan_dsp_2016}. The original lab has been converted from Matlab to Python and some of the contents has been modified. 
The lab demonstrates concepts from Chapters 2 and 3 in the text-book. The intention is to give a better understanding of how sinusoidal signals described by complex amplitude vectors, \emph{phasors}.

The first part of this lab gives training in sinusoidal signals and complex exponentials in Python. The last part uses this to estimate the direction of a sound source from the phase difference between signals received at two microphones. 
This type of direction estimates is the basis for steering and focusing ultrasound beams with \emph{phased arrays} in sonar and medical ultrasound. The principle is also used to steer radar beams and in wifi and 5G antennas.



\subsection{Software Tools: Python with Spyder and \jupyterlab}
The programs shall be written in Python. The module \numpy is used to represent signals and \matplotlib to plot graphs.
Including Python's \cmath and \rmath modules can make the code easier to read, but be aware that \cmath and \rmath only handle scalars. Note also that \cmath always returns a complex number, even if the imaginary part is zero. 
The recommended setup of modules for this lab is shown in \tableref{tab:import-libraries-code}

The recommended Python programming environment is \emph{Spyder}~\cite{raybaut_spyder_2024}, which is included in the \emph{Anaconda}~\cite{noauthor_anaconda_2024} package management. However, any Python package management and programming environment should function.

All code files shall be included with the lab report. We recommend collecting everything, text, code, and figures, into a single \jupyterlab Notebook~\cite{project_jupyter_jupyter_nodate}, but pdf with separate Python files is also acceptable.

\afterpage{\clearpage\input{../figs/table-import-libraries-lab1.tex}}

\section{Theory}

\subsection{Sinusoidal Signals and Phasors}
Recall from the lectures that a sinusoidal signal $x(t)$ is written as 
\begin{align}
	x(t)&= A \cos(\omega t + \phi) = \Re\left\{ A e^{j(\omega t + \phi)} \right\} 
		=  \Re\left\{ X e^{j\omega t} \right\} \quad , &
	X &= A e^{j\phi} 
\end{align}
where $A$ is the amplitude, $\omega = 2\pi f$ is the angular frequency, $f$ is the frequency, and $\phi$ is the phase. $X$ is a \emph{phasor} or \emph{complex amplitude} that includes both the amplitude and the phase of the signal.
Analysis of sinusoidal signals like $x(t)$ can be simplified by manipulating complex phasors  instead of working with the amplitude and phase separately, see Chapter 2 in the textbook \cite{mcclellan_dsp_2016} for details.


\subsection{Adding Sinusiods using Complex Exponentials}
Consider a signal that is the sum of sinusoids with equal frequency $f_0$ where the amplitudes $A_k$ and phases $\phi_k$ of the individual signals can be different, 
\begin{align}
	x_s(t)= \sum_{k=1}^{N} A_k \cos(2\pi f_0 t + \phi_k)   \:.
\end{align}
The resulting signal can be found by representing the individual components as complex vectors, summing these, and then taking the real part. This is called \emph{phasor summation} and is easier than using trigonometric identities, 
\begin{align}
	x_s(t)&= \Re\left\{ \sum_{k=1}^{N}  A_k e^{j\phi_k} e^{j2\pi f_0 t } \right\}  
			= \Re\left\{ \sum_{k=1}^{N}  X_k  e^{j2\pi f_0 t } \right\}  \qquad ,  &
	X_k &= A_k e^{j\phi_k} \:.
	\label{eq:phasorsum}
\end{align}
The factor $ e^{j2\pi f_0 t}$ in \eqref{eq:phasorsum} is equal for all the individual signals.
Hence, the amplitude $A_k$ and phase $\phi_k$ of the summed signal $x_s(t)$ can be found by summing the complex amplitudes, 
\begin{align}
	x_s(t)&= \Re \left\{ X_s e^{j\omega t } \right\} = A_s \cos(2\pi f_0 t + \phi_s)  \:,
	 &  X_s&= \sum_{k=1}^{N} X_k = A_s e^{j\phi_k}  \:.
\end{align}
The resulting signal will have the same frequency $f_0$ and period $T_0=1/f_0$ as the original signals. 

\subsection{Harmonics and Periodic signals}
Consider another signal $x_h(t)$ where the frequencies $f_k$ of the individual cosine-waves are different, but still integer multiples of a fundamental frequency $f_0$, 
\begin{align}
	f_k& = k f_0 \quad , & k&= 0,1,2, \ldots \:.
\end{align}
The individual signals $\cos(2\pi k f_0 t +\phi_k)$ are called \emph{harmonics}.
The summed signal $x_h(t)$ can be written as
\begin{align}
	x_h(t)= \sum_{k=1}^{N}  A_k \cos(2\pi k f_0 t +\phi_k)
	= \Re\left\{ \sum_{k=1}^{N}  X_k  e^{j2\pi k f_0 t } \right\}  \:.
	\label{eq:harmonicsum}
\end{align}
The period $T_0$ of the fundamental frequency $f_0$ is 
\begin{align}
	T_0 &= \frac{1}{f_0} \:,
\end{align}
while the periods $T_k$ of the harmonics are 
\begin{align}
	T_k &= \frac{1}{f_k} = \frac{1}{k f_0} = \frac{T_0}{k} \:.
\end{align}
The signal with frequency $kf_0$ will repeat itself after the period $T_k=T_0/k$, and will have repeated itself $k$ times after the fundamental period $T_0$. Hence, all the frequency components $kf_0$ will also be periodic with period $T_0$, and the resulting summed signal $x_h(t)$ will be periodic with period $T_0$ given by the fundamental frequency.

\section{Programming Tips}

\subsection{Complex Numbers in Python with \numpy}
Complex numbers are in Python treated like other numbers. The module \cmath contains elementary mathematical functions for use on \emph{scalar} complex numbers, while \numpy includes mathematical functions for use on \emph{arrays} of complex numbers.

\tableref{tab:complex-overview} lists the basic operations on complex numbers in Python.
A complex number in Python has three public members, \verb|real|, \verb|imag|, and \verb|conjugate()|. Other operations on can be found in \cmath or \numpy. The \numpy function calls in \tableref{tab:complex-overview} are prefixed by \verb|np| due to the way \numpy was imported, see \tableref{tab:import-libraries-code}, while the \cmath functions were imported individually. 

\input{../figs/table-complex-python}

\subsection{Displaying Phasors}
The Python file \verb|zplot.py| contains functions to plot complex numbers as phasors. 
This is done by the function \verb|zplot.phasor()|, see description in \tableref{tab:zplot}.

\input{../figs/table-zplot.tex}

\subsection{Vectorization}
The \numpy module allows mathematical operations to be used on arrays. This is convenient when defining signals as function of time. In the cosine-function $x(t)=A \cos(2\pi f t + \phi)$, the amplitude $A$ and phase $\phi$ are scalars, while the time $t$ is a vector spanning the time interval to be investigated. 
Vectors covering a defined interval can be created in in \numpy by one of the following methods.

\begin{enumerate}[1)]
	\item Specifying start, stop and step by the function \texttt{arange}. \\ 
	 A time vector $t$ is defined as 
	 \verb| t = np.arange(t_start, t_end, dt|), \\
	 where \verb|t_start| is the first point in the time-vector, \verb|t_end| marks the end of \verb|t|, and \verb|dt| is the interval between the time points. Note that the last value \verb|t_end| is not included in the time vector, \verb|t| ends on the last point before \verb|t_end|.
	 
	\item Specifying start, stop, and total number of points by the function \texttt{linspace}. \\
	A time vector $t$ is defined as \verb| t = np.linspace(t_start, t_end, n_points)|, \\
	where \verb|t_start| is the first point in \verb|t|, \verb|t_end| marks the end of \verb|t|, and \verb|n_points| is number of points in the vector. 
	Note again that the last value \verb|t_end| is not included.

	\item Specifying start, stop, and total number of points by the function \texttt{logspace}. \\ 
	This is the same as \verb|linspace|, but the numbers are evenly spaced on a logarithmic scale. The start and end points are specified by their logarithms, \verb|start=2| means the first value is \num{e2}=\num{100}.
\end{enumerate}
	
	
\subsection{Comparing graphs}
We sometimes need to compare graphs to see how similar they are.
Two useful methods to compare graphs in \matplotlib are
\begin{enumerate}[1)]
	\item Use \emph{subplots} to stack the graphs vertically or horizontally.
	
	\item Plot  the two curves in the same graph, the first with a solid line ('-') and the second with a dashed line ('- -'). Even if the curves are very similar, the first one will be visible behind the second.

\end{enumerate}

Example code is shown in \tableref{tab:subplot-example}, more information is found in the documentation for \matplotlib.
	
\input{../figs/table-subplot.tex}

\section{Training Exercises}

\subsection*{Reporting}
Collect answers and code in a \jupyterlab notebook. Export this to pdf and upload it to Canvas.

You may prefer to do some of the coding in another development tool, such as \spyder. The testing and debugging options are better in \spyder than in \jupyterlab. .py-files from from \spyder can be loaded into \jupyterlab and exported as separate pdf-files.

\subsection{Complex Numbers}

\begin{enumerate}[1)]
	\item Load \verb|zplot| and enter the two complex numbers $z_1 = 2e^{j\pi/3}$ and $z_2= -\sqrt{2} + 5j$. 
	
		Use Python to find the real and imaginary parts, magnitude, and phase of $z_1$ and $z_2$.
	
		Display $z_1$ and $z_2$ and the sum $z_1+z_2$ as phasors with \verb|zplot.phasor|. 		
		The input to \verb|zplot.phasor| is specified as a \verb|list|, this is made by enclosing the numbers in square brackets, e.g., \verb|[z1, z2]|.

	\item Find the complex conjugate $z^*$ and inverse $1/z$ for $z_1$ and $z_2$ and plot them together with $z_1$ and $z_2$ using \verb|zplot|.
	
		Recall what you have learned about complex numbers in math courses. Are the results as expected?
			
	\item Calculate the product $z_1 z_2$ and ratio $z_1/z_2$ and plot them using \verb|zplot|. 	
		Are these results as expected?

	\item Calculate the products of the conjugates, $z_1 z_1^*$ and $z_2 z_2^*$. 
	
	Plot them in the same diagram as $z_1$ and $z_2$ and explain the result.

	\item Calculate the sums $z_1+z_1^*$ and differences $z_1-z_1^*$ of the conjugates and plot them in the same diagram as $z_1$. Do the same for $z_2$. Explain these results.
		
\end{enumerate}
		
		
\subsection{Python Function to Generate a Sinusoid Signal }

\begin{enumerate}[1)]
	\item Write a function (\verb|def| in Python) that generates a single sinusoid, $x(t)= A \cos(2\pi f t + \phi)$ from the four input arguments amplitude $A$, frequency $f$, phase $\phi$ and duration. 
	
	The function shall return the signal $x(t)$ and the time vector $t$ where the signal is evaluated.
	
	The function shall generate exactly 32 values of the sinusoid per period. 

	A skeleton for this function \verb|make_cos| with the recommended function call and documentation string is listed in \tableref{tab:make-cos}. 

	\item Demonstrate that your function works by plotting the output for the following parameters:
		\begin{align*}
			A&=\num{e4} & f&= \qty{1.5}{MHz} &	\phi&=-\ang{45} & \text{Duration \qty{e-6}{s}}
		\end{align*}		
		Note that the phase must be converted to radians before calculating the result. 
		
		Calculate the value of $x(t)$ at $t=0$. Does this agree with the plot?
		
		What is the period of the signal? Does this agree with the plot?

\end{enumerate}

\input{../figs/table-make-cos-code}


\subsection{Python Function to Generate a Sum of Sinusoid Signals}
Signals are often described as a sum of sinusoids with different amplitudes, frequencies, and phases. It can therefore be convenient to have a function that generates a sum from several cosine-functions, each specified by its amplitude  $A_k$, frequency $f_k$, and phase $\phi_k$.

\begin{enumerate}[1)]
\item Write a function that generates a signal
	\begin{align}
		x(t)= \sum_{k=1}^{N} A_k \cos(2\pi f_k t + \phi_k) =\sum_{k=1}^{N} X_k e^{j2\pi f_k t}.
	\end{align}
	The input arguments to the function are the complex amplitude $X_k=A_k e^{j\phi}$, frequency $f_k$, sample rate $f_s$ and the signal duration.
	
	The function shall return the summed signal $x(t)$ and the time vector $t$ where the signal is evaluated.
	
	The frequencies $f_k$ and complex amplitudes $X_k$ shall be specified as lists or \numpy arrays, and the function shall accept any number of frequency components. 
	
	Each frequency $f_k$ shall match a complex amplitude $X_k$, so these vectors must have equal length. The resulting function must check for this and return an error message if the lengths are different. 
	
	A skeleton of this function \verb|make_summed_cos| with the recommended function call and documentation string is listed in \tableref{tab:summed_cos}. 
		
	\item Demonstrate that your function works by plotting the output for a signal that is the sum of the following components.
	\begin{center}
		\begin{tabular}{ccc}
				&	Frequency	& Complex amplitude	\\
			k	&	$f_k$ [Hz]		& $X_k$  			\\
			\hline
			1	&	\num{0}		& \num{10}				\\
			2	&	\num{100}	& $14e^{-j\pi/3}$		\\
			3	&	\num{250}	& 	$8j$				\\
		\hline
		\end{tabular}
	\end{center}
	Set the sample rate to \qty{10000}{Samples/s}, the duration of the signal to \qty{0.1}{s}, and the start time to \qty{0}{s}. Plot the result with \matplotlib.
	
	\item Measure the period $T_0 $ of the signal from the graph. Compare this to the periods $T_k$ of the individual frequency components $f_k$.
	
	Explain how the period of the summed signal can be calculated from the periods of the individual components.
	
\item Generate the signal
	\begin{align*}
		x(t) &= \Re{-2e^{j50\pi t} - e^{j50\pi(t-0.02)} +(2-3j)e^{j50\pi t}   } 
	\end{align*}
	over a time range that covers 2 periods. 
	
	Plot the signal $x$ a function of time $t$.
	
	\item All frequency components in the signal above are equal. Hence, the amplitude and phase can be calculated by summing its complex amplitudes, \emph{phasors}.
	
	Use the function \verb|zplot.phasor| from earlier to plot the phasor diagram for this signal, and check that this agrees with the result from \verb|make_summed_cos|
	
	\verb|zplot.phasor| has optional arguments that can be set to illustrate this better
	
	\begin{tabular}{ll}
		\verb|include_sum = True| & Include the sum of all the phasors to the plot. \\
		\verb|include_signal = True| & Plot the signals corresponding to the phasors. \\
		\verb|frequency = <value>|  & Frequency to use when plotting the signals.
		
	\end{tabular}

\end{enumerate}

\input{../figs/table-make-summed-cos-code.tex}

\section{Lab Exercise: Direction finding}
The text in this exercise is taken from \cite{mcclellan_lab_2016} and somewhat modified.

Why do humans have two ears? One answer is that the brain can process acoustic signals received at the two ears and determine the direction to the source of the acoustic energy. Using sinusoids, we can describe and analyze a simple scenario that explains this direction finding capability in terms of phase differences or time-delay differences. 
This principle is the basis for a wide range of other applications, such as  radars that locate and track airplanes, 5G and wifi antannas, and phased array transducers for medical ultrasound imaging and sonar.

\subsection*{Exercises: Direction Finding with Microphones}
Consider a system consisting of two microphones that both hear the same source signal. The microphones are placed some distance apart, so the sound must travel different paths from the source to the receivers. When the travel paths have different lengths, the two signals will arrive at different times. 
The time difference between the signals received by the two microphones allows us to calculate the direction to the source. If the source signal is a sinusoid, we can measure the travel time differences by measuring phases.
The scenario is illustrated in \figref{fig:overview}. A vehicle travelling along the roadway has a siren that  transmits a sinusoidal waveform with frequency $f_s$=\qty{400}{Hz}. The roadway forms the $x$-axis of a coordinate system.
The two microphones are located some distance away and aligned parallel to the roadway. The distance from the road to the microphones is $y_r$=\qty{100}{m} and the microphone separation is $d$=\qty{0.40}{m}. The task is to process the signals from the microphones to find the direction to the vehicle, described by the angle $\theta$ in \figref{fig:overview}.

\begin{figure}[!h]	
\begin{center}
	\begin{tikzpicture}[scale=3][>=stealth]

	\def\angle{33.7}
	\def\ymax{2.2}
	\def\xmax{2.0}
	\def\yr{1.8}
	\def\xv{1.2}
	\def\d{0.2}
	\def\sourcesize{0.03}
	
	% Colours
	% \colorlet{anglecolor}{blue} 
	\colorlet{microphonecolor}{blue!60!black}
	\colorlet{sourcecolor}{red!60!black}
	\colorlet{anglecolor}{sourcecolor} 
	\colorlet{roadcolor}{gray!20!black} 
	
	% Styles  
	\tikzstyle{axisline}=[very thick] 
	\tikzstyle{axisline}=[thick] 
	
	% Axes and road
	%\filldraw[fill=roadcolor, draw=black, nearly transparent] (-\xmax/2, -0.02) rectangle (\xmax*0.9, +0.02);
	\draw[axisline, ->] (-\xmax*0.7, 0) -- (\xmax, 0)  node[below]{$x$};
	\draw[axisline, ->] (0, -0.1) -- (0, \ymax)  node[right]{$y$};
	\filldraw[fill=sourcecolor, draw=black, opacity=0.7] (\xv-0.1, -0.05) rectangle  ++(0.2, 0.1) ;
	
	% Microphones
	\filldraw[fill=microphonecolor, draw=microphonecolor] (-\d/2, \yr) coordinate(M1) 
		circle (\sourcesize) 
		node[microphonecolor, above left]{M1};
	\filldraw[fill=microphonecolor, draw=microphonecolor] (\d/2, \yr) coordinate(M2)
		circle (\sourcesize) 		
		node[microphonecolor, above right]{M2};
	\draw[- ] (M1) -- node[above, midway, fill=white!] {$d$} (M2);

	% Sound paths
	\draw[sourcecolor, -] (\xv, 0) -- (M1) 
		node[midway, anchor=east]{$r_1$};
	\draw[sourcecolor, -] (\xv, 0) -- (M2) 
		node[midway, anchor=west]{$r_2$};;
	\draw[anglecolor, -, dashed] (\xv, 0) -- (0,\yr);
		
	% Angle
	\draw[draw=anglecolor] (0,\yr)++(0, -0.6)  arc(-90:-90+\angle:0.6) coordinate(ARC) 
		node[midway, below, anglecolor]{$\theta$};

	% Rulers
	\draw[|->, sourcecolor] (0, -0.15) -- ++(\xv, 0) node[at end, below]{$x_v$};
	\draw[|-|, microphonecolor] (-0.3, 0) -- ++(0, \yr) node[near end, left]{$y_r$};

\end{tikzpicture}


	\caption{Direction finding using two microphones. A vehicle at position $x_v$ travels along the $x$-axis while emitting a sound with frequency $f_s$=\qty{400}{Hz}. The sound is picked up by two microphones M1 and M2 positioned with spacing $d$=\qty{0.40}{m}. The difference in path length $\Delta r=r_1-r_2$ causes a phase-shift between the signals received by the two microphones. This phase shift can be used to estimate the direction to the vehicle, specified by the angle $\theta$. }
	\label{fig:overview}
\end{center}
\end{figure}

\begin{figure}[!h]	
\begin{center}
	\input{../figs/fig-direction-finding-approximation.tex}
	\caption{Zoomed-in version of \figref{fig:overview}. When the distance to the source is very long compared to the distance between the microphones, $r_1, r_2 \gg d$, the paths can be approximated as parallel and the difference in travel path length to the two microphones is is 
		$r_2 - r_1 = \Delta r\approx d \sin \theta$. }
	\label{fig:far_field_zoom}
\end{center}
\end{figure}


\begin{enumerate}[1)]
\item The time from the sound is transmitted by the source until it is received by one of the microphones can be computed for the two  propagation paths $r_1$ and $r_2$. 
	The time is given by the distance from the vehicle location at coordinate $(x_v,0)$ to either M1 at coordinate $(-\tfrac12 d, y_r)$ or M2 at $(+\tfrac12 d, y_r)$. 
	
	The speed of sound in air is $c$=\qty{340}{m/s}. 
	Write mathematical expressions for the time $t_1$ it takes for the sound to travel from the the source to M1 and for $t_2$ from the source to to M2. 
	Plot $t_1$ and $t_2$ as functions of the vehicle position $x_v$ from \qty{-400}{m} to \qty{+400}{m}.
	
	Can you see any difference between $t_1$ and $t_2$ in the graph?
	

\item In the simplest model, the signals received by the microphones, $s_1(t)$ at M1 and $s_2(t)$ at M2, are delayed copies of the transmitted signal $s(t)$,
	\begin{align}
		s_1(t)&= s(t-t_1)  &	s_2(t)&= s(t-t_2)  
	\end{align}
	where $s(t)$ is the signal transmitted from $x_v$.
	
	Assume that the source signal $s(t)$ is a zero-phase sinusoid at frequency $f$=\qty{400}{Hz} and set the amplitude of the transmitted signal to $A$=\num{1}. 	
	The phases $\phi_1$ and $\phi_2$ of the received signals $s_1$ and $s_2$ can be found from the delays $t_1$ and $t_2$. 
	
	Find the  phases $\phi_1$ and $\phi_2$. 
	
	Use the \verb|subplot| function in \verb|Matplotlib| and make a figure with two subplots. 
	
	Plot the phases $\phi_1$ and $\phi_2$ as functions of $x_v$ from \qty{-400}{m} to \qty{400}{m} in one subplot.
	
	Plot the phase difference $\Delta \phi = \phi_2 - \phi_1$ as function of $x_v$ in the second subplot.
	
	Comment the results.
	
\item The received signals can be represented as phasors. 
	
	Use \verb|zplot| to show the signals received by M1 and M2 when the vehicle is at positions $x_v$=\qty{-400}{m}$, \qty{-100}{m}$, \qty{0}{m}$, \qty{100}{m}$, and \qty{400}{m}.
	
	Set optional argument \verb|include_signal=True| to display the signals, and set the argument \verb|frequency| to the correct frequency. 
	
	Calculate the phase differences by hand and compare this with the plots.
	
	
\item How do we convert relative time-shift into the direction $\theta$?
	
	The distance from the microphones to the source is much larger than the distance between the microphones, making the the paths to M1 and M2 almost parallel. 
	This is illustrated in \figref{fig:far_field_zoom}, where we have zoomed in on the microphones in \figref{fig:overview}. 
	The difference $\Delta r$ in propagation distance for paths $r_1$ and $r_2$ can now approximated to
	\begin{align}
		\Delta r = r_1-r_2 \approx d \sin\theta \:.
		\label{eq:distance_difference}
	\end{align}
	This is called the \emph{far field approximation} and is often used to find the beam pattern from antennas, loudspeakers, ultrasound transducers, and other sources transmitting waves.
	
	Calculate the propagation time difference $\Delta t=t_2-t_1$ from the approximation \eqref{eq:distance_difference}. Use this to find the phase difference $\Delta \phi_F$ between the two received signals.
	Plot $\Delta \phi_F$ in the same graph as the  correct value $\Delta \phi$ found previously.
	
	Comment the result.
	

\item The objective is now to write a Python function that will process the received signals to find the direction $\theta$. 
	
	Show first that the phase difference between two phasors $X_1$ and $X_2$ is given by
	\begin{align*}
		\Delta \phi = -\angle \{ X_1 X_2^* \}
	\end{align*}
	where the superscript * denotes the complex conjugate. 
	
	Calculate the complex amplitudes $X_1 = A_1 e^{j\phi_1}$ and $X_2 = A_2 e^{j\phi_2}$ received at M1 and M2 as function of vehicle position $x_v$. Assume no loss, so that $A_1=A_2=A$.
	
	
	Use this and earlier results to write a Python-function that will compute the direction $\theta$ from the complex amplitudes.
		
	Run this function for the vehicle moving from $x_v=\qty{-400}{m}$ to \qty{+400}{m} and plot the angle $\theta_F$ calculated from the phase-shifts using the far-field approximation. Use degrees [\unit{\degree}] when plotting the angle.
	
	
	
	Compare $\theta_F$ to the true angle calculated from the actual position of the vehicle.
	
	Comment the result.
		
\end{enumerate}

\subsection*{Concluding remarks}
This exercise has illustrated how the direction to a sound source can be estimated using two receivers separated by a distance much smaller than the distance to the source.
This principle is the basis for how sound and radar beams can be steered and focused without moving parts in \emph{phased arrays}. Such arrays are used to steer medical ultrasound beams by GE Vingmed Ultrasound and in underwater sonars made by Kongsberg Discovery.


\bibliographystyle{ieeetr}
\bibliography{../defs/tse2280.bib}
\end{document}